%% CV for Alfa Laval - Backend Software Engineer (Dansk)
%% Compiled with: lualatex

\documentclass[11pt,a4paper,sans]{moderncv}
\moderncvstyle{banking}
\moderncvcolor{blue}

\renewcommand*{\firstnamestyle}[1]{{\fontsize{34}{36}\bfseries\upshape\color{color1}#1}}
\renewcommand*{\lastnamestyle}[1]{{\fontsize{34}{36}\bfseries\upshape\color{color1}#1}}
\renewcommand*{\sectionstyle}[1]{{\sectionfont\color{color1}#1}}

\usepackage[utf8]{inputenc}
\usepackage{hyperref}
\hypersetup{
    colorlinks=true,
    linkcolor=blue,
    filecolor=magenta,
    urlcolor=blue,
    pdftitle={Jacob El-Omar - CV},
    pdfpagemode=FullScreen,
}
\usepackage[scale=0.80]{geometry}
\usepackage{import}

\name{Jacob}{El-Omar}
\address{Aalborg, Danmark}{}{}
\phone[mobile]{+45 53 35 27 82}
\email{elomarjc@gmail.com}
\extrainfo{\href{https://linkedin.com/in/jacob-el-omar}{LinkedIn} \textbullet{} \href{https://github.com/elomarjc}{GitHub}}

\begin{document}

\makecvtitle

% ============================================================
%     PROFIL
% ============================================================

\vspace{6pt}
\small{Nyuddannet civilingeniør (cand.polyt.) i elektroniske systemer fra Aalborg Universitet med stærke kompetencer inden for C\#, .NET Core og relationelle databaser. Erfaren i at opbygge robuste backend-løsninger, optimere SQL-forespørgsler og strukturere versionsstyring med Git. Klar til at tage ejerskab over interne værktøjer og understøtte jeres digitale transformation i spændingsfeltet mellem klassisk softwareudvikling (IT) og industriel automation (OT).}

% ============================================================
%     FAGLIGE KOMPETENCER
% ============================================================

\section{Faglige Kompetencer}
\vspace{1pt}
\begin{itemize}

\item \textbf{Backend-udvikling \& C\#}: C\#, .NET Core, objektorienteret programmering (OOP), designmønstre (Dependency Injection, Singletons), REST API'er.

\item \textbf{Databaser \& SQL}: MSSQL, PostgreSQL, SQLite, database-migrationer, optimering af komplekse forespørgsler samt datamodellering.

\item \textbf{Automatisering \& OT-forståelse}: Programmering af indlejrede systemer i C/C++, PLC-forståelse, samt signalbehandling og reguleringsteknik.

\item \textbf{Versionsstyring \& CI/CD}: Erfaren med Git, GitHub Actions, samt akademisk kendskab til Azure DevOps og agile processer (Scrum).

\item \textbf{Tværfagligt samarbejde}: Uddannet via problembaseret læring (PBL) med fokus på tæt integration mellem softwarekrav og fysisk systemdesign.

\end{itemize}

% ============================================================
%     ERHVERVSERFARING
% ============================================================

\section{Erhvervserfaring}
\vspace{3pt}
\begin{itemize}

\item{\cventry{Jan 2026--Nu}{IT- \& Webudvikler (Frivillig)}{Dawah Danmark}{Aalborg, Danmark}{}{\vspace{1pt}
\begin{itemize}
    \item Udviklede og vedligeholdt interne systemer og databaser på WordPress-platformen.
    \item Optimerede database-arkitekturen, hvilket sikrede øget stabilitet og hurtigere indlæsningstider.
\end{itemize}}}

\vspace{3pt}

\item{\cventry{Sep 2024--Jun 2025}{AI/ML Engineer (Praktik)}{Ai Health Highway}{Remote}{}{\vspace{1pt}
\begin{itemize}
    \item Udviklede og testede algoritmer i Python til smart stethoskoper (AiSteth).
    \item Sikrede signalintegritet og dokumenterede kode i overensstemmelse med medicinske kvalitetsstandarder.
    \item Samarbejdede om integration af softwaren med indlejret hardware.
\end{itemize}}}

\vspace{3pt}

\item{\cventry{Jun 2023--Aug 2024}{Lager- og Logistikmedarbejder}{Lyn Express}{Aalborg, Danmark}{}{\vspace{1pt}
\begin{itemize}
    \item Styrede sortering og logistik af forsendelser under tidspres.
    \item Automatiserede administrative processer i Excel for at øge effektiviteten.
\end{itemize}}}

\end{itemize}

% ============================================================
%     UDDANNELSE
% ============================================================

\section{Uddannelse}
\vspace{1pt}
\begin{itemize}

\item{\cventry{Sep 2023--Jun 2025}{Civilingeniør, cand.polyt. (Elektroniske Systemer)}{Aalborg Universitet}{Aalborg, Denmark}{}{\vspace{1pt}
Speciale: ``Adaptive Noise Cancellation For Electronic Stethoscopes.'' Integrerede signalbehandling på realtidsplatforme.
}}

\vspace{3pt}

\item{\cventry{Sep 2020--Aug 2023}{Bachelor i teknisk videnskab (Elektronik og IT)}{Aalborg University}{Aalborg, Danmark}{}{\vspace{1pt}
Specialisering i reguleringsteknik og indlejrede systemer. Projekter inkluderede faldsikringssystemer på Cypress PSoC (C) og robotstyring.
}}

\end{itemize}

% ============================================================
%     UDVALGTE PROJEKTER
% ============================================================

\section{Udvalgte Projekter}
\vspace{1pt}
\begin{itemize}

\item{\textbf{Grow a Fish (Fritidsprojekt)}: Egenudviklet spil i Flutter/Dart. Implementerede en Supabase PostgreSQL-backend, REST-API'er, relationelt databasedesign og realtidssynkronisering.}

\item{\textbf{Embedded Faldsikringssystem (Universitetsprojekt, 4. sem.)}: Designede og programmerede et realtids fall-detection system på en Cypress PSoC mikrocontroller i C. Etablerede I2C-sensorkommunikation.}

\item{\textbf{Regulering af Traverskran (Bachelorprojekt, Universitetsprojekt)}: Udviklede en automatiseret model- og reguleringsarkitektur (PID/LQR) for traverskranstyring med feedbacksløjfer og systemidentifikation.}

\end{itemize}

% ============================================================
%     SPROG
% ============================================================

\section{Sprog}
\vspace{1pt}
\begin{itemize}
\item Dansk (modersmål), Engelsk (flydende).
\end{itemize}

% ============================================================
%     REFERENCER
% ============================================================

\section{References}
\vspace{1pt}
\begin{itemize}
\item{Afleveres efter aftale.}
\end{itemize}

\end{document}
